\section{Machine Learning: Extension}
\vspace{-0.5em}
\definition{Discriminative \& Generative Models}{Learning models can be categorized based on their conceptual approach to modeling data:
\begin{itemize}
\item \b{Discriminative Models:} Model only the decision function. Discrimination is often easier than full description (e.g., SVMs, Neural Networks).
\item \b{Generative Models:} Model the underlying concepts of the data and have the capability to generate new samples (e.g., Naive Bayes, HMM).
\end{itemize}
}
% \vspace{-0.5em}
\definition{Expected vs. Empirical Risk}{The formal goal of learning is risk minimization, but the true data distribution is usually unknown:
\begin{itemize}
\item \b{Expected Risk} \f{R_{\infty}(f_{\theta})}: Expected loss on data from a specific, unknown probability distribution \f{P(x,y)}.
\cf{R_{\infty}(f_{\theta})=\int L(f_{\theta}(x),y)dP(x,y)}
\item \b{Empirical Risk} \f{R_{n}(f_{\theta})}: The average loss calculated only on the \f{n} known training samples, which operates exclusively on the training data as a practical proxy.
\cf{R_{n}(f_{\theta})=\frac{1}{n}\sum_{i=1}^{n}L(f_{\theta}(x_{i}),y_{i})}
\end{itemize}
}
% \vspace{-0.5em}
\definition{Structural Risk Minimization (SRM)}{A framework that balances empirical risk with model complexity. This establishes a lower bound on the expected risk to prevent overfitting. It is achieved by adding a complexity penalty term \f{C(f_{\theta})} to the minimization problem, constrained by a regularization param \f{\lambda}.
\cf{g \mapsto argmin_{\theta \in \Theta} R_{n}(f_{\theta}) + \lambda C(f_{\theta})}
}
% \vspace{-0.5em}
\definition{Additional Learning Paradigms}{Beyond strict supervised and unsupervised learning, other specific paradigms include:
\begin{itemize}
\item \b{Semi-supervised:} Utilizes primarily unlabeled data alongside a small amount of labeled instances.
\item \b{Anomaly Detection:} An unsupervised technique focused on the detection of deviations from normality using generative or discriminative models. It outputs an anomaly score \f{Y=[0,1]} or distance \f{Y=\mathbb{R}^{+}} (e.g., One-class SVM).
\end{itemize}
}
% \vspace{-0.5em}
\definition{k-Nearest Neighbors (k-NN)}{A non-parametric supervised model that determines predictions based on the local neighborhood of data. During learning, it basically just stores the training data. The only parameter is the neighborhood size \f{k}, which acts as a form of regularization.}